\HistoricTitle{BEMERKUNGEN.}
{\small
\noindent\hspace*{\parindent}Die vorstehenden von \textsc{Gauss} zwischen dem 21. Oktober 1810 und dem 5. November 1812 geschriebenen Briefstellen gewähren einen Einblick in die verschiedenartigen Gedankenverbindungen, die \textsc{Gauss} um die Zeit der Abfassung der \textit{Disquisitiones circa seriem} an die Reihe \(F(\alpha,\beta,\gamma,x)\) angeknüpft hat. Das Interesse von \textsc{Gauss} am Integrallogarithmus geht, wie der Werke II, S. 444 abgedruckte Brief an \textsc{Encke} zeigt, bis in die Göttinger Studentenzeit zurück\textsuperscript{\(\dagger\)}), dem oben abgedruckten Briefe [5.] an \textsc{Laplace} entnehmen wir, daß \textsc{Gauss} sich auch schon 1800 mit der inexplikabeln Funktion \(\Psi(x)\) beschäftigt hat.

\indent Zu der Stelle des Briefes [3] an \textsc{Bessel}, wo von der \textls[100]{Fortsetzung} der Reihe \(F(\alpha,\beta,\gamma,x)\) auf Grund eines »höheren Prinzips« die Rede ist, vergleiche man die Schlußsätze des nachgelassenen art. 38. der Abhandlung \textit{circa seriem}, Werke III, S. 209–210. Die bei \textsc{Acunha} (1790) auftretende Erklärung des Logarithmus durch Umkehrung der Exponentialreihe, die \textsc{Gauss} in demselben Briefe lobend erwähnt, findet sich
\par\medskip\noindent\rule{.24\linewidth}{.4pt}\par\smallskip
\noindent\hspace*{\parindent}\textsuperscript{\(\dagger\)}) Vergl. auch den Abschnitt [I.] des folgenden Hauptstücks \textit{Zur Lehre von den Reihen}.
\par\smallskip\noindent\hfill48\textsuperscript{*}
\clearpage\SourceLeaf{387}\ArithmeticHead{380}{L}{ANALYSIS BRIEFWECHSEL.}
\noindent auch in den beiden Preisschriften, die \textsc{Wolfgang} und \textsc{Johann Bolyai} 1837 der \textsc{Jablonowski}schen Gesellschaft der Wissenschaften eingereicht hatten; siehe P. \textsc{Stäckel}, \textit{Wolfgang und Johann Bolyai, Geometrische Untersuchungen}, Leipzig 1913, I, S. 129, 245 und II, S. 228. Über die Bezeichnung \(\sin^2\varphi,\cos^2\varphi\) hat sich \textsc{Gauss} in ähnlicher Weise wie hier S. 365 im Jahre 1814 auch öffentlich geäußert, siehe Werke IV, S. 361; vergl. dazu auch \textit{Briefwechsel zwischen Gauss und Schumacher} III, S. 292.

\indent Zu den auf die Integration im komplexen Gebiet bezüglichen Erörterungen des Briefes [4.] an \textsc{Bessel}\textsuperscript{*}) vergleiche man das im folgenden als Abschnitt [VI.] der \textit{Lehre von den Reihen} abgedruckte Bruchstück \textit{Bestimmung der Convergenz der Reihen u.s.w.}

\indent Über die von \textsc{Gauss} in dem Briefe [5.] berührte Aufgabe der Wahrscheinlichkeitsrechnung vergleiche man die Bemerkung zu der Aufzeichnung Nr. 113 in dem unten folgenden Abdruck des \textit{Tagebuchs}. Zu der in demselben Briefe erörterten Prioritätsfrage, oben S. 373–374, ist noch folgendes zu bemerken. In einem an \textsc{Gauss} gerichteten Briefe\textsuperscript{**}) vom 31. Mai 1809 schreibt \textsc{Legendre} in Bezug auf die Methode der kleinsten Quadrate:

\indent »Je ne vous dissimulerai-donc pas, Monsieur, que j’ai éprouvé quelque regret de voir qu’en citant mon mémoire pag. 221[\textsuperscript{***})], vous dites \textls[100]{principium nostrum quo jam inde ab anno 1795 usi sumus} etc. Il n’est aucune découverte qu’on ne puisse s’attribuer en disant qu’on avoit trouvé la même chose quelques années auparavant; mais si on n’en fournit pas la preuve en citant le lieu où on l’a publiée, cette assertion devient sans objet et n’est plus qu’une chose désobligeante pour le véritable auteur de la découverte. En Mathématiques il arrive très souvent qu’on trouve les mêmes choses qui ont été trouvées par d’autres et qui sont bien connues; c’est ce qui m’est arrivé nombre de fois, mais je n’en ai point fait mention et je n’ai jamais appelé \textls[100]{principium nostrum} un pr[incipe] qu’un autre avait publié avant moi. Vous êtes assez riche de [votre] fonds, Monsieur, pour n’avoir rien à envier à personne; et [je suis] bien persuadé au reste que j’ai à me plaindre de l’expression seulement et nullement de l’intention \({.}\;{.}\;{.}\;{.}\;{.}\;{.}\;{.}\)«

\indent Daß \textsc{Gauss} in seinem Briefe auf diese Prioritätsfrage einging, war durch einen an ihn gerichteten Brief von \textsc{Laplace} vom 15. November 1811 verursacht, in dem es in bezug auf die Methode der kleinsten Quadrate heißt:

\indent »Monsieur \textsc{gauss} dit dans son ouvrage sur le mouvement elliptique qu’il la possedait avant que M. \textsc{le gendre} l’ait publiée; je desirerois bien sçavoir si avant cette publication on avait imprimé quelque chose en allemagne sur cette méthode et je prie Monsieur \textsc{gauss} de vouloir bien m’en instruire.«

\indent In bezug auf den \textls[100]{Auszug} aus der Abhandlung über die Anziehung der elliptischen Sphäroide, den \textsc{Gauss} dem Briefe [7.] beigelegt hatte (siehe oben S. 379) vergleiche man die Briefe an \textsc{Schumacher} vom 31. Dezember 1812, \textit{Briefwechsel zwischen Gauss und Schumacher} I, 1860, S. 95, und an \textsc{Olbers} vom 8. April 1813, \textit{Wilhelm Olbers, Sein Leben und seine Werke} II, 1, 1900, S. 515. In seinem Antwortschreiben vom 20. November 1812 schreibt \textsc{Laplace} darüber das folgende:

\indent »A votre lettre était jointe une demonstration très simple du theorème sur les attractions des spheroïdes elliptiques. j’ai commencé par le lire et j’ai été ravi de sa simplicité. elle a produit le même effet sur ceux des membres de l’institut qui se sont occupés de cet objet et aux quels je l’ai communiquée suivant vos intentions. cette demonstration a quelque analogie avec une demonstration fort simple du même
\par\medskip\noindent\rule{.24\linewidth}{.4pt}\par\smallskip
\noindent\hspace*{\parindent}\textsuperscript{*}) Dieser Teil des Briefes ist auch schon Werke VIII, S. 90–92 abgedruckt.
\par\noindent\hspace*{\parindent}\textsuperscript{**}) Der Brief befindet sich im \textsc{Gauss}archiv; die in der folgenden Wiedergabe eingeklammerten Worte sind aus der Handschrift beim Öffnen des Briefes weggeschnitten worden.
\par\noindent\hspace*{\parindent}\textsuperscript{***}) Werke VII, 1906, S. 253.
\clearpage\SourceLeaf{388}\ArithmeticHead{381}{R}{BEMERKUNGEN.}
\noindent theorème, donnée par M. \textsc{jvori}, geometre anglais, dans la seconde partie des transactions philosophiques de 1809[\textsuperscript{*})]. peut être, vous ne la connaissés pas; je vous engage à la voir. \({.}\;{.}\;{.}\;{.}\)«

\indent Diese Mitteilung von \textsc{Laplace} erwähnt \textsc{Gauss} in dem \textit{Additamentum} zu seiner Abhandlung, Werke V, S. 21 und in seiner \textit{Anzeige} vom 5. April 1813 in den Göttingischen Gelehrten Anzeigen, Werke V, S. 279, siehe insbesondere S. 285.

\indent Die vorstehend wiedergegebenen Auszüge aus Briefen von \textsc{Laplace} sind ebenso wie die Briefstellen [1.], [3.], [4.], [6.] an \textsc{Bessel} nach den im \textsc{Gauss}archiv befindlichen Handschriften, die an \textsc{Laplace} nach den ebenda befindlichen Abschriften abgedruckt. Es sind Abschriften von 8 Briefen von \textsc{Gauss} an \textsc{Laplace} vorhanden, deren Urschriften sich in Paris im Privatbesitz befinden, die Abschriften sind 1877 angefertigt und von J. \textsc{Bertrand}, dem damaligen vorsitzenden Sekretär der Pariser Akademie der Wissenschaften, beglaubigt.
\begin{flushright}\textsc{Schlesinger}.\end{flushright}
\par\medskip\noindent\rule{.24\linewidth}{.4pt}\par\smallskip
\noindent\hspace*{\parindent}\textsuperscript{*}) \textsc{James Ivory}, \textit{On the Attractions of homogeneous Ellipsoids} Philosophical Transactions of the Royal Society of London 1809, S. 345.
\par
\begin{center}\rule{.30\linewidth}{.4pt}\end{center}
\clearpage}
